% =============================================================================
%  HBUE-demo.tex — 湖北经济学院 Beamer 模板示例
%  Hubei University of Economics Beamer Template Demo
% =============================================================================
%  编译: xelatex HBUE-demo.tex (推荐 XeLaTeX, 需编译两遍以生成进度条)
%
%  主题选项:
%    \usetheme{HBUE}         默认: 校徽红 + 浅色
%    \usetheme[blue]{HBUE}   官网蓝主色调
%    \usetheme[dark]{HBUE}   深色模式
%    \usetheme[nologo]{HBUE} 不加载校徽
% =============================================================================

\documentclass[aspectratio=169, 10pt, UTF8]{ctexbeamer}

\usetheme{HBUE}

% ── 元信息 ──
\title[HBUE Beamer 模板]{湖北经济学院\\Beamer 演示文稿模板}
\subtitle{基于 \LaTeX{} Beamer 的学术演示解决方案}
\author[张三]{张\ 三}
\supervisor{李四\ 教授}
\institute[HBUE]{湖北经济学院\ \ 信息管理学院\\Hubei University of Economics}
\date{\today}

% ── 可选包 ──
\usepackage{booktabs}
\usepackage{amsmath,amssymb}

\begin{document}

% ═══════════════════════════════════════════
%  标题页
% ═══════════════════════════════════════════
{
\setbeamertemplate{headline}{}
\setbeamertemplate{footline}{}
\begin{frame}[plain, noframenumbering]
  \titlepage
\end{frame}
}

% ═══════════════════════════════════════════
%  目录
% ═══════════════════════════════════════════
\tocpage

% ═══════════════════════════════════════════
%  1. 模板简介
% ═══════════════════════════════════════════
\section{模板简介}

\begin{frame}{设计理念}
  \begin{columns}[T]
    \begin{column}{0.52\textwidth}
      \textbf{设计参考:}
      \begin{itemize}
        \item SJTUBeamer — 上海交通大学
        \item USTCBeamer — 中国科学技术大学
        \item University-Presentation-Beamer
        \item Metropolis 现代极简风格
      \end{itemize}
      \vspace{0.8em}
      \textbf{技术栈:}
      \begin{itemize}
        \item XeLaTeX + ctexbeamer
        \item TikZ 矢量绘图
        \item 16:9 宽屏比例
      \end{itemize}
    \end{column}
    \begin{column}{0.48\textwidth}
      \begin{block}{视觉识别 (VIS)}
        校训: 厚德博学\ 经世济民\\
        校徽: 腾飞的经院 (凤凰 + E + U)\\
        校徽红: \texttt{\#CA001D} (主色)\\
        凤凰黄: \texttt{\#FFF100} (点缀)\\
        官网蓝: \texttt{\#084B9D} (可选主色)
      \end{block}
      \begin{alertblock}{编译要求}
        请使用 \textbf{XeLaTeX} 编译\\
        Overleaf 需切换编译器
      \end{alertblock}
    \end{column}
  \end{columns}
\end{frame}

\begin{frame}{主题选项}
  本模板支持以下主题选项组合:
  \vspace{0.5em}
  \begin{columns}[T]
    \begin{column}{0.48\textwidth}
      \begin{exampleblock}{校徽红主题 (默认)}
        \texttt{\textbackslash usetheme\{HBUE\}}
      \end{exampleblock}
      \begin{exampleblock}{深色模式}
        \texttt{\textbackslash usetheme[dark]\{HBUE\}}
      \end{exampleblock}
    \end{column}
    \begin{column}{0.48\textwidth}
      \begin{exampleblock}{官网蓝主题}
        \texttt{\textbackslash usetheme[blue]\{HBUE\}}
      \end{exampleblock}
      \begin{exampleblock}{隐藏校徽}
        \texttt{\textbackslash usetheme[nologo]\{HBUE\}}
      \end{exampleblock}
    \end{column}
  \end{columns}
  \vspace{0.5em}
  选项可自由组合,如 \texttt{[blue, dark]}.
\end{frame}

% ═══════════════════════════════════════════
%  2. 基础排版
% ═══════════════════════════════════════════
\section{基础排版}

\begin{frame}{列表与强调}
  \begin{columns}[T]
    \begin{column}{0.48\textwidth}
      \textbf{无序列表:}
      \begin{itemize}
        \item 第一层级
        \begin{itemize}
          \item 第二层级
          \begin{itemize}
            \item 第三层级
          \end{itemize}
        \end{itemize}
        \item 回到第一层级
      \end{itemize}
      \vspace{0.5em}
      \textbf{有序列表:}
      \begin{enumerate}
        \item 第一步
        \item 第二步
        \item 第三步
      \end{enumerate}
    \end{column}
    \begin{column}{0.48\textwidth}
      \textbf{文字强调:}
      \begin{itemize}
        \item \textbf{粗体强调}
        \item {\kaishu 楷体强调}
        \item \alert{警告强调}
        \item \hl{主色高亮}
        \item \hlg{金色高亮}
      \end{itemize}
      \vspace{0.5em}
      \textbf{描述列表:}
      \begin{description}
        \item[\LaTeX] 排版系统
        \item[Beamer] 演示文档类
      \end{description}
    \end{column}
  \end{columns}
\end{frame}

\begin{frame}{块环境 (Blocks)}
  \begin{block}{普通信息块}
    这是一个普通的 \textbf{block} 环境,用于展示一般性内容。
    支持多行文本,自动换行。
  \end{block}

  \begin{alertblock}{提醒块}
    这是一个 \textbf{alertblock},用于需要特别注意的信息。
  \end{alertblock}

  \begin{exampleblock}{示例块}
    这是一个 \textbf{exampleblock},用于展示示例或最佳实践。
  \end{exampleblock}
\end{frame}

% ═══════════════════════════════════════════
%  3. 数学公式
% ═══════════════════════════════════════════
\section{数学公式}

\begin{frame}{公式排版}
  \textbf{行内公式:}
  欧拉恒等式 $\mathrm{e}^{i\pi} + 1 = 0$ 被称作最美的数学公式。

  \vspace{0.4em}
  \textbf{行间公式:}
  \begin{equation}
    \int_{-\infty}^{+\infty} \mathrm{e}^{-x^{2}} \, dx = \sqrt{\pi}
  \end{equation}

  \begin{theorem}[拉格朗日中值定理]
    若 $f \in C[a,b] \cap D(a,b)$, 则 $\exists\,\xi \in (a,b)$ 使得
    \begin{equation*}
      f'(\xi) = \frac{f(b) - f(a)}{b - a}.
    \end{equation*}
  \end{theorem}
\end{frame}

\begin{frame}{矩阵与多行公式}
  \textbf{矩阵:}
  \begin{equation}
    \mathbf{A} =
    \begin{bmatrix}
      a_{11} & a_{12} & \cdots & a_{1n} \\
      a_{21} & a_{22} & \cdots & a_{2n} \\
      \vdots & \vdots & \ddots & \vdots \\
      a_{m1} & a_{m2} & \cdots & a_{mn}
    \end{bmatrix}
  \end{equation}

  \textbf{多行对齐:}
  \begin{align}
    \nabla \cdot \mathbf{E}   &= \frac{\rho}{\varepsilon_0} \\
    \nabla \cdot \mathbf{B}   &= 0 \\
    \nabla \times \mathbf{E}  &= -\frac{\partial \mathbf{B}}{\partial t}
  \end{align}
\end{frame}

% ═══════════════════════════════════════════
%  4. 表格与图片
% ═══════════════════════════════════════════
\section{表格与图片}

\begin{frame}{三线表}
  \begin{table}
    \centering
    \caption{湖北经济学院办学历史沿革}
    \vskip0.5em
    \begin{tabular}{ccl}
      \toprule
      \textbf{序号} & \textbf{时间} & \textbf{重大事项} \\
      \midrule
      1 & 1907 年 & 张之洞创办"湖北商业中学堂"(办学源头) \\
      2 & 2002 年 & 三校合并组建湖北经济学院 \\
      3 & 2012 年 & 开始硕士研究生教育 \\
      4 & 2021 年 & 获批硕士学位授予单位 \\
      \bottomrule
    \end{tabular}
  \end{table}
\end{frame}

\begin{frame}{图片插入}
  \begin{figure}
    \centering
    \begin{beamercolorbox}[wd=0.55\textwidth, center, sep=0.9cm,
      rounded=true]{block body}
      \Large\color{HBUEgray}%
      图片占位区域\\
      \vskip0.4em
      \footnotesize
      \texttt{\textbackslash includegraphics}\\
      \texttt{[width=\textbackslash textwidth]}\\
      \texttt{\{your-figure.pdf\}}
    \end{beamercolorbox}
    \caption{图片插入示例 (使用占位框)}
  \end{figure}
\end{frame}

% ═══════════════════════════════════════════
%  5. 多栏与动画
% ═══════════════════════════════════════════
\section{高级排版}

\begin{frame}{多栏布局}
  \begin{columns}[T]
    \begin{column}{0.48\textwidth}
      \begin{block}{左栏要点}
        \texttt{columns} 环境支持:
        \begin{itemize}
          \item 2--4 栏自由分割
          \item 顶部/居中/底部对齐
          \item 百分比宽度控制
        \end{itemize}
        非常适合对比展示。
      \end{block}
    \end{column}
    \begin{column}{0.48\textwidth}
      \begin{block}{右栏要点}
        常用比例分配:
        \begin{itemize}
          \item $0.5 + 0.5$ 平分
          \item $0.6 + 0.4$ 黄金比
          \item $0.3 + 0.7$ 三七开
        \end{itemize}
        总宽度不超过 1.0 即可。
      \end{block}
    \end{column}
  \end{columns}

  \vspace{0.5em}
  \begin{quote}
    多栏布局适合信息对比、图文并茂、中英对照等场景。
  \end{quote}
\end{frame}

\begin{frame}{逐步显示}
  \begin{itemize}
    \item<1-> 第一项 — 从第 1 步开始可见
    \item<2-> 第二项 — 从第 2 步开始可见
    \item<3-> 第三项 — 从第 3 步开始可见
    \item<4-> 第四项 — 从第 4 步开始可见
  \end{itemize}

  \vspace{0.8em}
  \visible<5->{%
    \begin{alertblock}{提示}
      使用 \texttt{<n->} 语法或 \texttt{\textbackslash pause}
      可实现逐条逐步显示。
    \end{alertblock}
  }
\end{frame}

% ═══════════════════════════════════════════
%  6. 引用与参考
% ═══════════════════════════════════════════
\section{参考文献}

\begin{frame}[allowframebreaks]{参考文献}
  \begin{thebibliography}{99}
    \bibitem{sjtubeamer}
    SJTUG.
    \newblock SJTUBeamer: 上海交通大学 Beamer 模板.
    \newblock GitHub, 2024.

    \bibitem{ustcbeamer}
    USTC TeX Users Group.
    \newblock USTCBeamer: 中国科学技术大学 Beamer 模板.
    \newblock GitHub, 2024.

    \bibitem{lamport}
    Leslie Lamport.
    \newblock \LaTeX: A Document Preparation System.
    \newblock Addison-Wesley, 2nd edition, 1994.

    \bibitem{beamer}
    Till Tantau, Joseph Wright, Vedran Mileti\'{c}.
    \newblock The Beamer Class User Guide.
    \newblock CTAN, 2024.
  \end{thebibliography}
\end{frame}

% ═══════════════════════════════════════════
%  致谢
% ═══════════════════════════════════════════
\thankyoupage

\end{document}
